\section{第二天}
\begin{enumerate}
  \item 今天是两个脑膜炎级别考点，
        也就是哪怕背成脑膜炎也要背下来的知识点，
        重要性程度直接无死角拉满！
        第一个考点，
        尽管在过去的试题当中已经有过考查，
        但基于考点的性质和功能，
        完全无法排除它作为常驻考点的可能性，
        务必重视！
        为避免影响发送，
        敏感词语作了替代，
        应该不影响阅读理解。
  \item \bs{法治思想}的概念和价值。
        \begin{enumerate}
          \item 法治思想是什么？
                \begin{enumerate}
                  \item 法治思想是\zx{着眼于}中华民族伟大复兴\zx{战略全局}和当今世界\zx{百年未有}之大变局，
                        \zx{顺应}实现中华民族伟大复兴\zx{时代要求}应运而生的\bs{重大战略思想}，
                  \item 也是\bs{新时代}中国特色社会主义思想的\bs{重要组成部分}。
                  \item 法治思想\zx{深刻回答}了\bs{新时代}为什么要实行全面依法治国、怎样实行全面依法治国等一系列重大问题，
                        饱含\bs{原创性、系统性、时代性、人民性和实践性}的理论品格，
                  \item \bs{为}\zx{深入推进}全面\zx{依法治国}、\zx{加快建设}社会主义\zx{法治国家}，运用\zx{制度威力}应对\zx{风险挑战}，
                  \item \zx{实现}党和国家长治久安，全面建设\zx{社会主义现代化国家}、实现中华民族伟大复兴的中国梦，
                  \item 提供了\bs{科学指南}。
                \end{enumerate}
          \item 法治思想的价值意义。
                \begin{enumerate}
                  \item \blh 是\zx{马克思主义法治理论}同中国\zx{法治建设}具体实际相结合、同\zx{中华优秀传统法律文化}相结合的\bs{最新成果}。
                  \item \blh 是对党领导法治建设丰富\zx{实践}和宝贵\zx{经验}的\bs{科学总结}。
                  \item \blh 是在\zx{法治轨道}上全面建设社会主义现代化国家的\bs{根本遵循}。
                  \item \blh 是引领法治中国建设，实现高质量\zx{发展}的\bs{思想旗帜}。
                \end{enumerate}
        \end{enumerate}
  \item 全面依法治国之政治方向一：\bs{坚持党的领导}。
        \begin{enumerate}
          \item 坚持党的领导的概念和价值。
                \begin{enumerate}
                  \item 党的领导是\zd{全面依法治国的法治}之魂，也是中国特色社会主义法治和西方资本主义法治最大的区别。
                  \item 党政军民学，东西南北中，党是\zx{领导一切}的。
                  \item 中国共产党是中国特色社会主义事业的\zx{坚强领导核心}，是\zx{最高政治领导力量}，
                  \item 各个领域、各个方面都必须\zx{坚定自觉坚持}党的领导。
                  \item 党的领导是中国特色社会主义\zx{最本质的特征}，是\zd{社会主义法治}\zx{最根本的保证}。
                  \item 坚持党的领导。
                        \begin{enumerate}
                          \item 是\zd{社会主义法治}的\zx{根本要求}，
                          \item 是党和国家的\zx{根本所在、命脉所在}，
                          \item 是全国各族人民的\zx{利益所系、幸福所系}，
                          \item 是\zd{题中/材料中/xx 所指出}的\zd{全面推进依法治国}的\zx{应有之义}。
                        \end{enumerate}
                  \item \zx{党的领导}和\zd{社会主义法治}\zx{是一致}的，\zd{社会主义法治}\zx{必须坚持}党的领导，党的领导\zx{必须依靠}\zd{社会主义法治}。
                  \item \zd{全面推进依法治国，建设社会主义法治国家}，只有在党的领导下才能\zx{有目的、有步骤、有秩序}地进行。
                \end{enumerate}
          \item 如何坚持党的领导（小标题必背，后面的阐述部分，可以照猫画虎，自己发挥）。
                \begin{enumerate}
                  \item 全面依法治国就是要\zx{坚持和改进}党的领导。
                        \begin{enumerate}
                          \item \zd{全面依法治国}是一项\zx{系统的工程}，
                          \item 只有\zx{发挥党}总揽全局、协调各方的\zx{领导核心作用}，
                          \item 才能完成\zd{全面依法治国}这一“系统工程”的\zx{总规划}，才能实现\zd{全面依法治国}的\zx{总目标}。
                        \end{enumerate}
                  \item 把党的领导\zx{贯彻落实到}依法治国的\zx{全过程和各方面}。
                        \begin{enumerate}
                          \item 必须坚持党领导立法、保证执法、支持司法、带头守法，
                          \item 把依法治国基本方略\zx{同依法执政基本方式统一}起来，
                          \item 充分发挥党的领导对\zd{全面依法治国}各个方面工作的\zx{保障作用}。
                        \end{enumerate}
                  \item 坚持把\zx{党的领导、人民当家作主、依法治国统一}起来（这个部分，要求完整背诵！可能单独命题）。
                        \begin{enumerate}
                          \item 三统一中\zx{最根本的是}坚持党的领导。党的领导是人民当家作主和\zd{依法治国}的\zx{根本保证}。
                          \item 人民当家作主是\bs{社会主义民主政治}的\zx{本质特征}，
                          \item \zd{依法治国}是党领导人民治理国家的\zx{基本方式}，
                          \item 三者统一是我国\bs{社会主义民主政治}\zx{伟大实践}。
                          \item \zd{人民代表大会制度}。
                                \begin{itemize}
                                  \item 是坚持党的领导、人民当家作主、\zd{依法治国}有机统一的\zx{根本制度安排}。
                                  \item 是实现党的\zx{领导和执政}的制度\zx{载体和依托}，是人民当家作主的\zx{根本途径和实现形式}。
                                \end{itemize}
                        \end{enumerate}
                  \item \zx{健全党领导}\zd{全面依法治国}的\zx{制度和工作机制}。\bs{最后都可以回到国家治理上}。
                        \begin{enumerate}
                          \item 党委、政法委要\zx{带头}在\zd{宪法法律}\zx{范围内活动}，
                          \item \zx{善于运用}\zd{法治}\zx{思维}和\zd{法治}\zx{方式，领导}\zd{政法}\zx{工作}，
                          \item 在推进国家\zx{治理体系和治理能力现代化}中发挥重要作用。
                        \end{enumerate}
                \end{enumerate}
        \end{enumerate}
\end{enumerate}
